使用 \cite{r-75}, Theorem 4.10 与 \cite{r-30}, Section 2.2 中的射影换元。只需作有限次旋转，即可将方向上确界限制到 \(|\omega_n|\ge c_n\)，其中 \(c_n>0\) 取某个固定值。将任一单位立方体 \(Q\) 平移到原点，并记 \(x=(x',t_0)\in Q\)。当 \(\rho\) 充分小时，限制后取上确界所用的每根管均位于 \(5Q\) 内，落在薄层 \(c_n/4\le |s-t_0|\le2\) 中，且 \(s-t_0\) 的符号固定。固定高度 \(t_0\) 和其中一个符号。在这一有界薄层及其像上，映射
\[
 (y',s)\longmapsto (Y',Y_n)=\left(y'/(s-t_0),\,1/(s-t_0)\right)
\]
与它的逆映射均一致 Lipschitz，且 Jacobian 行列式的绝对值有统一正下界和上界。对任意这类管，其轴为 \(y'=x'+(s-t_0)v\)，变换后的直线为 \(Y'=v+Y_nx'\)，方向为 \((x',1)/|(x',1)|\)，与 \(v\) 无关。轴线段的像长度有界，管的像包含于轴线段像的 \(C_n\rho\) 邻域。故用有界数目的同方向单位管即可覆盖此像，这些管的半径为 \(C_n\rho\)。原管以及每根覆盖管的体积均与 \(\rho^{n-1}\) 可比。

将 \(|f|\) 限制到 \(5Q\) 与固定符号薄层的交集上，再经射影映射输送，并在其余地方延拓为零。所得 \(F\) 依赖于 \(t_0\) 和所取的符号，适用于该切片中的所有 \(x'\)。由 Jacobian 界，其 \(L^n\) 范数为 \(\lesssim\|f\|_{L^n(5Q)}\)。比较测度可知，每个原来的平均值均由 \(K_{C_n\rho}F((x',1)/|(x',1)|)\) 的常数倍控制。在该立方体内有界的 \(x'\) 范围上，方向映射为单射，面积 Jacobian 有正下界。应用假设的 Kakeya 估计，得到切片的 \(L^n\) 范数由 \(C L(1/\rho)^{2A}\|f\|_{L^n(5Q)}\) 控制。对两个符号求和并对 \(t_0\) 积分。然后在各单位格点立方体上平移，利用其伸缩体 \(5Q\) 的重叠次数有统一上界。当 \(\rho\) 有固定正下界时，各管包含在以 \(x\) 为中心的固定大小的球内，且各管体积有正下界，故 Young 不等式直接给出所述估计。
